\paragraph{Definition 3.5 (Time-Bounded Computational Depth [9,35]).}
For time bounds $s\in\mathbb N$ and $t\in\mathbb N\cup\{\infty\}$ and a string $x\in\{0,1\}^*$, the $(s,t)$-time-bounded computational depth of $x$ is defined as
\[\operatorname{cd}^{s,t}(x):=K^s(x)-K^t(x).\]
We omit the superscript $t$ if $t=\infty$.
\paragraph{Definition 6.1.}
For a polynomial $\tau:\mathbb N^3\to\mathbb N$, we define the promise problem $\mathsf{Gap}_{\tau}\mathsf{MINcKT}$ to be $(\Pi_{\mathrm{Yes}},\Pi_{\mathrm{No}})$ such that
\[\Pi_{\mathrm{Yes}}:=\{(x,y,1^t,1^s):K^t(x\mid y)\leq s-\operatorname{cd}^{t,\tau(|x|,|y|,t)}(y)\},\]
\[\Pi_{\mathrm{No}}:=\{(x,y,1^t,1^s):K^{\tau(|x|,|y|,t)}(x\mid y)>s+\log\tau(|x|,|y|,t)\}.\]
\paragraph{Definition 6.5.}
For a polynomial $\tau:\mathbb N^3\to\mathbb N$ and a function $\sigma:\mathbb N\to\mathbb N$, we define the promise problem $\mathsf{Gap}_{\tau,\sigma}\mathsf{MINcKT}$ to be $(\Pi_{\mathrm{Yes}},\Pi_{\mathrm{No}})$ such that
\[\Pi_{\mathrm{Yes}}:=\{(x,y,1^t,1^s):K^t(x\mid y)\leq s-\operatorname{cd}^{t,\tau(|x|,|y|,t)}(y)\},\]
\[\Pi_{\mathrm{No}}:=\{(x,y,1^t,1^s):K^{\tau(|x|,|y|,t)}(x\mid y)>\sigma(|x|)\cdot s+\log\tau(|x|,|y|,t)\}.\]
Note that $\mathsf{Gap}_{\tau,1}\mathsf{MINcKT}=\mathsf{Gap}_{\tau}\mathsf{MINcKT}$ is trivially reducible to $\mathsf{Gap}_{\tau,\sigma}\mathsf{MINcKT}$ for every $\sigma\geq1$.
\paragraph{Definition 6.7 (Access Structure).}
An access structure $\mathcal A\subseteq2^{[n]}$ is a ``monotone'' collection of subsets of $[n]$; that is, for every $T\supseteq S\in\mathcal A$, we have $T\in\mathcal A$. The minimum weight of $\mathcal A$ is defined to be $w(\mathcal A):=\min\{|T|:T\in\mathcal A\}$.
We will prove that there is a generic reduction from the problem of estimating the weight of access structures $\mathcal A$ to $\mathsf{GapMINcKT}$ if there exists an efficient secret sharing scheme for $\mathcal A$.
\paragraph{Definition 6.8 (Secret Sharing [12]).}
A secret sharing scheme for $\mathcal A$ is a pair $(\mathsf{Share},\mathsf{Rec})$ of a randomized algorithm $\mathsf{Share}$ and a deterministic algorithm $\mathsf{Rec}$ with the following properties for every $\ell\in\mathbb N$:
\begin{enumerate}
\item Correctness: For every $T\in\mathcal A$ and for every string $x\in\{0,1\}^\ell$, any output of $\mathsf{Share}(x)$ is a sequence $(y_1,\ldots,y_n)$ of $n$ strings that satisfies $\mathsf{Rec}(y_T)=x$, where $y_T:=\{(i,y_i):i\in T\}$.
\item Privacy: For every $T\notin\mathcal A$ and for every random variable $X$ on $\{0,1\}^\ell$, the random variables $X$ and $\mathsf{Share}(X)_T$ are statistically independent.
\end{enumerate}
We observe that the privacy condition can be stated in terms of Kolmogorov complexity. We mention in passing that Kolmogorov complexity-theoretic versions of privacy conditions are studied in, e.g., [10,48].
\paragraph{Definition 6.10.}
A family $\mathcal A=\{\mathcal A_\phi\}_{\phi\in\{0,1\}^*}$ of access structures is said to admit efficient secret sharing schemes if there exists a pair $(\mathsf{Share},\mathsf{Rec})$ of a randomized polynomial-time algorithm $\mathsf{Share}$ and a deterministic polynomial-time algorithm $\mathsf{Rec}$ such that for every $\phi\in\{0,1\}^*$, the pair $(\mathsf{Share}(\phi,-),\mathsf{Rec}(\phi,-))$ is a secret sharing scheme for the access structure $\mathcal A_\phi$.
Benaloh and Leichter [13] showed that access structures represented by monotone formulas admit efficient secret sharing schemes.
\paragraph{Lemma 6.11 ([13]).}
Let $\mathcal A:=\{\mathcal A_\phi\}_{\phi\in\{0,1\}^*}$ be the family of access structures $\mathcal A_\phi:=\{T\subseteq[n]:\phi(\chi_T)=1\}$, where $\phi$ is a monotone formula on $n$ variables and $\chi_T\in\{0,1\}^n$ denotes the characteristic vector of $T\subseteq[n]$. Then, $\mathcal A$ admits efficient secret sharing schemes.
\paragraph{Definition 6.12 (Minimum Monotone Satisfying Assignment; MMSA).}
For a monotone formula $\phi$ on $n$ variables, the weight of an assignment $\alpha\in\{0,1\}^n$ is defined to be $\sum_{i=1}^n\alpha_i$. Let $\mathsf{MMSA}(\phi)$ denote the minimum weight of $\alpha\in\{0,1\}^n$ such that $\phi(\alpha)=1$.
Observe that $\mathsf{MMSA}(\phi)=w(\mathcal A_\phi)$ for the family $\mathcal A$ of access structures of Lemma 6.11.
